\section{Structural amortization as a cost-to-capability hypothesis}

Reuse is not free. Let $C_I$ be the one-time cost of inducing, validating and storing a reusable structure. Let $c_B$ be the baseline per-task cost of solving without reuse and $c_R$ the per-task cost of retrieving, adapting and verifying the reusable structure. Then
\[
C_B(n)=n c_B,
\qquad
C_R(n)=C_I+n c_R.
\]
When $c_R<c_B$, reuse becomes strictly cheaper once
\[
n>\frac{C_I}{c_B-c_R}.
\]
The executable reference implementation reports this break-even count and returns no finite break-even when the reuse path is not cheaper per task. This elementary accounting is important because prior workflow-compilation work can reduce runtime calls without establishing net efficiency when offline compilation cost is omitted \cite{elyadouni2026tracecompiler}.

The empirical paper uses a fuller vector
\[
C_{\mathrm{total}}=
C_{\mathrm{induction}}+
C_{\mathrm{training}}+
C_{\mathrm{retrieval}}+
C_{\mathrm{adapt/reason}}+
C_{\mathrm{tools}}+
C_{\mathrm{verification}}.
\]
No coordinate is silently set to zero. A cost-to-capability point is admissible only if it reaches the frozen capability target without crossing the registered validity-failure ceiling.

\begin{hypothesis}[Structural transfer]
At matched semantic similarity, a witnessed structural score provides material incremental prediction of held-out transfer success.
\end{hypothesis}

\begin{hypothesis}[Semantic-decoy rejection]
A witness-aware transfer gate has a lower invalid-transfer rate than semantic retrieval or analogy-only baselines on Q3 cases with high semantic and low structural similarity.
\end{hypothesis}

\begin{hypothesis}[Static structural training efficiency]
At a frozen structural-OOD capability target, static structural curation reaches the target with fewer training tokens or examples than strong semantic/diversity and skill-graph parents in regimes containing cross-domain structural redundancy.
\end{hypothesis}

\begin{hypothesis}[Inference efficiency]
For a frozen base model and target capability, reuse of validated structural objects/operators reduces test-time reasoning/search cost relative to semantic memory, analogy prompting and strong skill/workflow-prior parents after all retrieval, adaptation and verification costs are included.
\end{hypothesis}

\begin{hypothesis}[Redundancy mechanism]
Under a controlled data generator, the advantage of structural reuse increases with structural redundancy after surface-semantic diversity is held approximately fixed.
\end{hypothesis}

\begin{hypothesis}[Shared substrate]
A structure identifier induced or registered during training/data curation can be reused at inference without changing to an unrelated task-specific representation. The strongest paper claim requires this shared object; two independent tricks would support two narrower results rather than one unification.
\end{hypothesis}

\subsection{Optional training-time extension: learner-conditioned structural saturation}
\label{sec:learner-conditioned-saturation}

The static training hypothesis above does not by itself imply that a structural family should receive the same training allocation throughout learning. Recent work already establishes substantial prior art for model-aware or online data selection, including evolving influence-based selection, online distribution optimization, model-state-dependent perplexity schedules and explicit missing-skill targeting \cite{yu2024mates,jiang2025ado,elhattami2025sst,he2026stat}. The present paper therefore claims no novelty for adaptive curricula, model-aware selection, saturation-aware reweighting, skill graphs or missing-skill profiles in general.

A narrower RAKL extension remains open because the transfer substrate is not a flat skill label. For a registered structural object $S$, a future learner-conditioned training view may estimate a scoped mastery vector
\[
m_t(S)=\bigl(m_P,m_C,m_B,m_R,m_T,m_F\bigr),
\]
where the coordinates denote principle acquisition, composition mastery, boundary robustness, representation invariance, cross-domain transfer and retention/forgetting resistance. The vector is learner-specific: it is bound to a model/checkpoint and a frozen probe family, and it can become stale after later parameter updates. It is not scientific authority.

The load-bearing distinction is that one coordinate may saturate while another remains unsaturated. In particular,
\[
\operatorname{mastered}(A)+\operatorname{mastered}(B)
\not\Rightarrow
\operatorname{mastered}(A+B),
\]
and high in-domain principle accuracy does not establish boundary or cross-domain transfer mastery. This prevents ``already knows the skill'' from becoming a scalar excuse to remove all further examples containing that principle.

The training-time extension should therefore be interpreted as budget reallocation rather than destructive deduplication. Repeated exposure may still be useful for robustness, calibration, representation invariance and forgetting prevention; prior LLM curation work has shown that intelligently repeated data can outperform random single-use exposure \cite{tirumala2023d4}. A learner-conditioned scheduler must preserve a registered repetition/retention floor and preferentially move only the \emph{marginal} budget toward unsaturated compositions, boundaries, representations, domains or hostile near-misses.

\begin{hypothesis}[Learner-conditioned structural allocation --- extension hypothesis]
In a training distribution with surface-diverse structural redundancy, an adaptive scheduler conditioned on a frozen learner-specific mastery vector over registered relational structures reaches a preregistered structural-OOD capability target at lower total cost than \emph{static RAKL structural curation} and the strongest faithful model-aware/skill-aware parent, while satisfying registered composition, boundary, retention and negative-transfer constraints.
\end{hypothesis}

The critical comparison is therefore
\[
\boxed{\text{Adaptive RAKL}-\text{Static RAKL}},
\]
not merely RAKL versus random sampling. If this contrast is measured but indistinguishable, the learner-conditioned extension is not established even if static structural curation remains useful. If a simpler parent such as STAT- or MATES-style selection matches the result, the claim must narrow to any residual value contributed by the directional structural/boundary/shared-substrate machinery. No result for this extension is claimed in the present manuscript.

Each hypothesis has a direct falsifier. Flat structural coefficients, high Q3 transfer, no break-even cost, a strong MASS/Skill-It/MATES/STAT/SWIFT-like parent matching RAKL, adaptive RAKL failing to beat static RAKL, harmful forgetting/composition trade-offs, or incompatible training/inference representations all narrow the paper.
